\subsection{The Boolean Representation Architecture (\texttt{bool})}

Python's \texttt{bool} type represents truth values. It has only two ordinary values:

\begin{verbatim}
True
False
\end{verbatim}

These are not strings. They are not numbers written with special spelling. They are boolean objects.

\begin{verbatim}
is_ready = True
is_finished = False

print(f"is_ready:    {is_ready}")
print(f"is_finished: {is_finished}")

print()
print(f"type(is_ready):    {type(is_ready)}")
print(f"type(is_finished): {type(is_finished)}")

print()
print("selected bool methods:")
print(f"{'__and__':<12} {'__and__' in dir(is_ready)}")
print(f"{'__or__':<12} {'__or__' in dir(is_ready)}")
print(f"{'__invert__':<12} {'__invert__' in dir(is_ready)}")
print(f"{'bit_length':<12} {'bit_length' in dir(is_ready)}")
print(f"{'to_bytes':<12} {'to_bytes' in dir(is_ready)}")
\end{verbatim}

Possible output:

\begin{verbatim}
is_ready:    True
is_finished: False

type(is_ready):    <class 'bool'>
type(is_finished): <class 'bool'>

selected bool methods:
__and__      True
__or__       True
__invert__   True
bit_length   True
to_bytes     True
\end{verbatim}

The presence of methods such as \texttt{bit\_length()} and \texttt{to\_bytes()} is not accidental. In Python, \texttt{bool} is closely connected to \texttt{int}. That connection is useful, but it also creates traps.

\subsubsection{Boolean Values as Singleton Objects: \texttt{True} and \texttt{False}}

There is only one normal \texttt{True} object and one normal \texttt{False} object in a Python process. They are singleton objects.

This means that every ordinary use of \texttt{True} refers to the same object, and every ordinary use of \texttt{False} refers to the same object.

\begin{verbatim}
a = True
b = True

c = False
d = False

print(f"a == b: {a == b}")
print(f"a is b: {a is b}")

print()
print(f"c == d: {c == d}")
print(f"c is d: {c is d}")

print()
print(f"id(a): {id(a)}")
print(f"id(b): {id(b)}")
print(f"id(c): {id(c)}")
print(f"id(d): {id(d)}")
\end{verbatim}

Possible output:

\begin{verbatim}
a == b: True
a is b: True

c == d: True
c is d: True

id(a): 140734431495200
id(b): 140734431495200
id(c): 140734431495232
id(d): 140734431495232
\end{verbatim}

The exact \texttt{id()} values do not matter. The important part is that \texttt{a is b} and \texttt{c is d} are both true.

Conceptually:

\begin{verbatim}
a ----\
      ---> True singleton object
b ----/

c ----\
      ---> False singleton object
d ----/
\end{verbatim}

The constructor \texttt{bool()} also returns one of these two singleton objects:

\begin{verbatim}
x = bool(42)
y = bool("No soup for you.")
z = bool("")

print(f"x: {x}, x is True:  {x is True}")
print(f"y: {y}, y is True:  {y is True}")
print(f"z: {z}, z is False: {z is False}")

print()
print(f"type(x): {type(x)}")
print(f"type(y): {type(y)}")
print(f"type(z): {type(z)}")
\end{verbatim}

Possible output:

\begin{verbatim}
x: True, x is True:  True
y: True, y is True:  True
z: False, z is False: True

type(x): <class 'bool'>
type(y): <class 'bool'>
type(z): <class 'bool'>
\end{verbatim}

The function \texttt{bool()} does not create a new private truth object. It decides whether the input is truthy or falsy and returns either \texttt{True} or \texttt{False}.

Some common truth tests:

\begin{verbatim}
values = [0, 1, 42, "", "D'oh!", [], [1, 2, 3], None]

for value in values:
    truth_value = bool(value)
    print(f"{repr(value):<12} -> {truth_value:<5} type={type(truth_value)}")
\end{verbatim}

Possible output:

\begin{verbatim}
0            -> False type=<class 'bool'>
1            -> True  type=<class 'bool'>
42           -> True  type=<class 'bool'>
''           -> False type=<class 'bool'>
"D'oh!"      -> True  type=<class 'bool'>
[]           -> False type=<class 'bool'>
[1, 2, 3]    -> True  type=<class 'bool'>
None         -> False type=<class 'bool'>
\end{verbatim}

This introduces truth-value testing. It does not mean that all these objects become booleans. The object \texttt{42} remains an integer. The object \texttt{"D'oh!"} remains a string. The call \texttt{bool(...)} produces a boolean result.

\subsubsection{Boolean as a Subclass of Integer: Numeric Compatibility and Practical Warnings}

In Python, \texttt{bool} is a subclass of \texttt{int}. This means that \texttt{True} and \texttt{False} can participate in many integer operations.

\begin{verbatim}
print(f"isinstance(True, bool): {isinstance(True, bool)}")
print(f"isinstance(True, int):  {isinstance(True, int)}")

print()
print(f"True == 1:  {True == 1}")
print(f"False == 0: {False == 0}")

print()
print(f"int(True):  {int(True)}")
print(f"int(False): {int(False)}")
\end{verbatim}

Possible output:

\begin{verbatim}
isinstance(True, bool): True
isinstance(True, int):  True

True == 1:  True
False == 0: True

int(True):  1
int(False): 0
\end{verbatim}

This explains why boolean values have some integer-like methods:

\begin{verbatim}
value = True

print(f"value: {value}")
print(f"type:  {type(value)}")

print()
print(f"value.bit_length(): {value.bit_length()}")
print(f"value.bit_count():  {value.bit_count()}")

print()
print(f"True + True:   {True + True}")
print(f"True + False:  {True + False}")
print(f"False + False: {False + False}")
\end{verbatim}

Possible output:

\begin{verbatim}
value: True
type:  <class 'bool'>

value.bit_length(): 1
value.bit_count():  1

True + True:   2
True + False:  1
False + False: 0
\end{verbatim}

This can be useful when counting true conditions:

\begin{verbatim}
has_ticket = True
is_late = False
knows_python = True

true_count = has_ticket + is_late + knows_python

print(f"true_count: {true_count}")
print(f"type(true_count): {type(true_count)}")
\end{verbatim}

Possible output:

\begin{verbatim}
true_count: 2
type(true_count): <class 'int'>
\end{verbatim}

The result is an integer because arithmetic with booleans uses their integer compatibility.

However, this compatibility should not be abused. Booleans are meant to express truth, not ordinary quantities.

Bad style:

\begin{verbatim}
is_valid = True
total = is_valid + 41

print(total)
\end{verbatim}

Possible output:

\begin{verbatim}
42
\end{verbatim}

This works, but it is unclear. A reader must remember that \texttt{True} behaves like \texttt{1}. If the program means a number, use a number. If the program means a truth value, use a boolean.

Better:

\begin{verbatim}
is_valid = True

if is_valid:
    total = 42
else:
    total = 0

print(total)
\end{verbatim}

The practical rule is:

\begin{quote}
\texttt{True} is numerically compatible with \texttt{1}, and \texttt{False} is numerically compatible with \texttt{0}, but boolean names should normally describe conditions, not quantities.
\end{quote}

Use names like:

\begin{verbatim}
is_ready = True
has_access = False
can_continue = True
found_match = False
\end{verbatim}

These names tell the reader that the value is meant to answer a yes/no question.

\subsubsection{Logical Operators vs. Bitwise Operators: \texttt{and} / \texttt{or} / \texttt{not} Compared with \texttt{\&} / \texttt{|} / \texttt{\textasciitilde{}}}

Python has logical operators and bitwise operators. They are not the same thing.

The logical operators are:

\begin{verbatim}
and
or
not
\end{verbatim}

They are used for truth-value logic.

\begin{verbatim}
has_ticket = True
is_late = False

print(f"has_ticket and not is_late: {has_ticket and not is_late}")
print(f"has_ticket or is_late:      {has_ticket or is_late}")
print(f"not has_ticket:             {not has_ticket}")
\end{verbatim}

Possible output:

\begin{verbatim}
has_ticket and not is_late: True
has_ticket or is_late:      True
not has_ticket:             False
\end{verbatim}

The bitwise operators are:

\begin{verbatim}
&
|
~
\end{verbatim}

For integers, they operate on binary bits:

\begin{verbatim}
a = 0b1010
b = 0b1100

print(f"a:     {a:04b}")
print(f"b:     {b:04b}")

print()
print(f"a & b: {a & b:04b}")
print(f"a | b: {a | b:04b}")
print(f"~a:    {~a}")
\end{verbatim}

Possible output:

\begin{verbatim}
a:     1010
b:     1100

a & b: 1000
a | b: 1110
~a:    -11
\end{verbatim}

The result of \texttt{\textasciitilde{}a} may look surprising because Python integers are not fixed-width unsigned buffers. Conceptually, \texttt{\textasciitilde{}x} is equal to \texttt{-x - 1}.

\begin{verbatim}
x = 42

print(f"x:     {x}")
print(f"~x:    {~x}")
print(f"-x - 1: {-x - 1}")
print(f"same:  {~x == -x - 1}")
\end{verbatim}

Possible output:

\begin{verbatim}
x:     42
~x:    -43
-x - 1: -43
same:  True
\end{verbatim}

Because booleans are integer-compatible, bitwise operators also work on booleans:

\begin{verbatim}
print(f"True & True:   {True & True}")
print(f"True & False:  {True & False}")
print(f"True | False:  {True | False}")
print(f"False | False: {False | False}")
\end{verbatim}

Possible output:

\begin{verbatim}
True & True:   True
True & False:  False
True | False:  True
False | False: False
\end{verbatim}

This may look the same as logical behavior for simple boolean values:

\begin{verbatim}
print(f"True and False: {True and False}")
print(f"True & False:   {True & False}")

print()
print(f"True or False:  {True or False}")
print(f"True | False:   {True | False}")
\end{verbatim}

Possible output:

\begin{verbatim}
True and False: False
True & False:   False

True or False:  True
True | False:   True
\end{verbatim}

But the operators are still different.

First difference: \texttt{and} and \texttt{or} do not necessarily return a boolean. They return one of the original operands.

\begin{verbatim}
first = "No soup for you."
second = "D'oh!"

print(f"first and second: {first and second!r}")
print(f"first or second:  {first or second!r}")

print()
print(f"type(first and second): {type(first and second)}")
print(f"type(first or second):  {type(first or second)}")
\end{verbatim}

Possible output:

\begin{verbatim}
first and second: "D'oh!"
first or second:  'No soup for you.'

type(first and second): <class 'str'>
type(first or second):  <class 'str'>
\end{verbatim}

Second difference: \texttt{and} and \texttt{or} short-circuit. They may skip the right side when the result is already decided.

\begin{verbatim}
print("case 1")
result1 = False and (1 / 0)
print(f"result1: {result1}")

print()
print("case 2")
result2 = True or (1 / 0)
print(f"result2: {result2}")
\end{verbatim}

Possible output:

\begin{verbatim}
case 1
result1: False

case 2
result2: True
\end{verbatim}

No division error occurs because the right side is not evaluated.

Bitwise operators do not short-circuit in the same way:

\begin{verbatim}
try:
    result = False & (1 / 0)
except ZeroDivisionError as err:
    print("operation failed")
    print(f"message: {err}")
    print(f"type(err): {type(err)}")
\end{verbatim}

Possible output:

\begin{verbatim}
operation failed
message: division by zero
type(err): <class 'ZeroDivisionError'>
\end{verbatim}

The expression on the right side was evaluated before the bitwise operation could be applied.

Third difference: \texttt{not} and \texttt{\textasciitilde{}} are completely different.

\begin{verbatim}
print(f"not True:  {not True}")
print(f"not False: {not False}")

print()
print(f"~True:  {~True}")
print(f"~False: {~False}")
\end{verbatim}

Possible output:

\begin{verbatim}
not True:  False
not False: True

~True:  -2
~False: -1
\end{verbatim}

Why?

Because \texttt{True} is numerically compatible with \texttt{1}, and \texttt{False} is numerically compatible with \texttt{0}.

\begin{verbatim}
print(f"~1: {~1}")
print(f"~0: {~0}")
\end{verbatim}

Possible output:

\begin{verbatim}
~1: -2
~0: -1
\end{verbatim}

So \texttt{\textasciitilde{}} is not a logical not operator. It is bitwise inversion.

The practical rule is:

\begin{quote}
Use \texttt{and}, \texttt{or}, and \texttt{not} for logical conditions. Use \texttt{\&}, \texttt{|}, and \texttt{\textasciitilde{}} only when you intentionally want bit-level operations or when a specific library type defines those operators for its own purpose.
\end{quote}

A final comparison:

\begin{verbatim}
is_python = True
is_clear = True

print("logical condition")
print(is_python and is_clear)

flags_a = 0b1010
flags_b = 0b1100

print()
print("bitwise flags")
print(f"{flags_a & flags_b:04b}")
\end{verbatim}

Possible output:

\begin{verbatim}
logical condition
True

bitwise flags
1000
\end{verbatim}

The practical summary is:

\begin{quote}
Python booleans are singleton truth-value objects. They are also integer-compatible, with \texttt{True} behaving numerically like \texttt{1} and \texttt{False} like \texttt{0}. This compatibility explains some behavior, but ordinary logical code should use \texttt{and}, \texttt{or}, and \texttt{not}, not the bitwise operators.
\end{quote}